% ============================================================
%  ANALYSE DE DONNEES DE SANTE AVEC PYTHON
%  Notes de cours -- VERSION CONDENSEE 6 HEURES
%  Compilation : xelatex -shell-escape notes_6h.tex (2 passes)
% ============================================================

\documentclass[11pt,a4paper]{article}

% --- Encodage et langue ---
\usepackage{fontspec}
\usepackage{polyglossia}
\setmainlanguage{french}
\setotherlanguage{english}

% --- Mise en page ---
\usepackage[margin=2.4cm,top=2.8cm,bottom=2.8cm]{geometry}
\usepackage{parskip}
\setlength{\parindent}{0pt}
\setlength{\parskip}{0.7em}

% --- Headers and footers ---
\usepackage{fancyhdr}
\pagestyle{fancy}
\fancyhf{}
\fancyhead[L]{\small\itshape Analyse de donn\'ees avec Python pour professionnels de la sant\'e}
\fancyhead[R]{\small\itshape Version condens\'ee 6 heures}
\fancyfoot[C]{\thepage}
\setlength{\headheight}{14pt}
\renewcommand{\headrulewidth}{0.3pt}

% --- Math ---
\usepackage{amsmath,amssymb,amsthm,mathtools}

% --- Quotes francaises + citations bibliographiques ---
\usepackage{csquotes}
\usepackage[round,authoryear]{natbib}

% Polyglossia ne definit pas \og / \fg (contrairement a babel-french).
% On les definit pour rester compatible avec la syntaxe habituelle.
\providecommand{\og}{\guillemotleft\,}
\providecommand{\fg}{\,\guillemotright}

% --- Tables ---
\usepackage{booktabs,array,longtable}

% --- Couleurs ---
\usepackage[svgnames,dvipsnames]{xcolor}
\definecolor{primary}{HTML}{0072B2}
\definecolor{primarydark}{HTML}{004F7A}
\definecolor{accent}{HTML}{009E73}
\definecolor{warn}{HTML}{D55E00}
\definecolor{codebg}{HTML}{F5F7F9}

% --- Code (minted) ---
\usepackage{minted}
\setminted{
  fontsize=\small,
  bgcolor=codebg,
  frame=none,
  framesep=2mm,
  breaklines=true,
  breakanywhere=true,
  python3=true,
  tabsize=2,
}
\setminted[python]{python3=true}

% --- Callout boxes (tcolorbox) ---
\usepackage{tcolorbox}
\tcbuselibrary{theorems,breakable,skins}

\newtcolorbox{astuce}{
  enhanced,colback=accent!4,colframe=accent!75,
  left=8pt,right=8pt,top=6pt,bottom=6pt,
  boxrule=0pt,leftrule=3pt,arc=2pt,
  title={\bfseries ASTUCE},
  fonttitle=\small\bfseries\sffamily,
  attach boxed title to top left={xshift=0pt,yshift=-3pt},
  boxed title style={colback=accent,colframe=accent,sharp corners},
  breakable
}

\newtcolorbox{attention}{
  enhanced,colback=warn!5,colframe=warn!75,
  left=8pt,right=8pt,top=6pt,bottom=6pt,
  boxrule=0pt,leftrule=3pt,arc=2pt,
  title={\bfseries ATTENTION},
  fonttitle=\small\bfseries\sffamily,
  attach boxed title to top left={xshift=0pt,yshift=-3pt},
  boxed title style={colback=warn,colframe=warn,sharp corners},
  breakable
}

\newtcolorbox{notebox}{
  enhanced,colback=primary!4,colframe=primary!75,
  left=8pt,right=8pt,top=6pt,bottom=6pt,
  boxrule=0pt,leftrule=3pt,arc=2pt,
  title={\bfseries NOTE},
  fonttitle=\small\bfseries\sffamily,
  attach boxed title to top left={xshift=0pt,yshift=-3pt},
  boxed title style={colback=primary,colframe=primary,sharp corners},
  breakable
}

\newtcolorbox{clinique}{
  enhanced,colback=accent!4,colframe=accent!60,
  left=8pt,right=8pt,top=6pt,bottom=6pt,
  boxrule=0pt,leftrule=3pt,arc=2pt,
  title={\bfseries CONTEXTE CLINIQUE},
  fonttitle=\small\bfseries\sffamily,
  attach boxed title to top left={xshift=0pt,yshift=-3pt},
  boxed title style={colback=accent!80!black,colframe=accent!80!black,sharp corners},
  breakable
}

% --- Section formatting ---
\usepackage{titlesec}
\titleformat{\section}
  {\Large\bfseries\color{primarydark}}
  {\thesection.}{0.6em}{}
\titleformat{\subsection}
  {\large\bfseries\color{primary}}
  {\thesubsection}{0.6em}{}
\titlespacing*{\section}{0pt}{1.8em}{0.8em}
\titlespacing*{\subsection}{0pt}{1.4em}{0.5em}

% --- Liens ---
\usepackage{hyperref}
\hypersetup{
  colorlinks=true,
  linkcolor=primarydark,
  citecolor=accent,
  urlcolor=primarydark,
  pdfauthor={Notes de cours},
  pdftitle={Analyse de donn\'ees avec Python pour professionnels de la sant\'e -- Version 6h}
}

% --- TOC formatting ---
\usepackage{tocloft}
\renewcommand{\cftsecfont}{\bfseries}
\renewcommand{\cftsecpagefont}{\bfseries\color{primary}}

\begin{document}

% ============================================================
% PAGE DE GARDE
% ============================================================

\thispagestyle{empty}
\vspace*{2cm}

\begin{center}

% --- Banner-style header block ---
\begin{tcolorbox}[
  enhanced,
  colback=primary,
  colframe=primary,
  boxrule=0pt,
  arc=8pt,
  left=22pt,right=22pt,top=20pt,bottom=20pt,
  width=\linewidth,
]
\color{white}
\begin{flushleft}
{\small\sffamily\bfseries\uppercase{Notes de cours}\hfill\uppercase{Version condens\'ee 6 heures}}\\[12pt]
{\Huge\bfseries Analyse de Donn\'ees avec Python}\\[4pt]
{\LARGE\bfseries pour Professionnels de la Sant\'e}\\[18pt]
{\large\itshape Six modules d'une heure -- du Python de base au machine learning clinique}
\end{flushleft}
\end{tcolorbox}

\vspace{2.2cm}

\begin{tabular}{rl}
\bfseries Public cibl\'e   & Professionnels de sant\'e, \'epid\'emiologistes, \\
                           & \'etudiants en sant\'e publique \\[6pt]
\bfseries Pr\'erequis      & Aucun en programmation. Quelques notions \\
                           & de statistique de base sont utiles. \\[6pt]
\bfseries Format           & Notes accompagnant 6 notebooks Jupyter \\[6pt]
\bfseries Langue           & Fran\c{c}ais \\[6pt]
\bfseries Outils           & Python 3.10+, Jupyter / Google Colab \\
\end{tabular}

\vfill

{\small\color{gray!70}
Mat\'eriel p\'edagogique libre. Si vous le r\'eutilisez, citez la source.
}

\end{center}

\newpage

% ============================================================
% TABLE DES MATIERES
% ============================================================

\tableofcontents
\thispagestyle{empty}

\newpage

% ============================================================
% AVANT-PROPOS
% ============================================================

\section*{Avant-propos}
\addcontentsline{toc}{section}{Avant-propos}

\subsection*{Ce que ces notes sont -- et ne sont pas}

Ces notes accompagnent les six notebooks Jupyter de la \textbf{version
condens\'ee 6 heures} du cours \emph{Analyse de donn\'ees de sant\'e
avec Python}. Elles d\'eveloppent la mati\`ere th\'eorique et
m\'ethodologique. Les notebooks contiennent le code complet et
ex\'ecutable ; ces notes contiennent les \textbf{concepts}, les
\textbf{formules}, les \textbf{contextes d'usage} et les
\textbf{pi\`eges} -- avec seulement quelques courts extraits de code
pour illustrer.

Si vous lisez ces pages \emph{avant} la s\'eance, vous saurez de quoi
on va parler et pourquoi. Si vous les relisez \emph{apr\`es}, elles
servent d'aide-m\'emoire compact.

\subsection*{Pourquoi 6 heures}

Un cours d\'etaill\'e de 30 heures existe (dossier \texttt{code/notebooks/}
du d\'ep\^ot). Mais pour une journ\'ee de formation, un atelier de trois
demi-journ\'ees, ou un cours d'introduction, six heures sont souvent ce
qu'on peut donner -- et ce qu'un apprenant peut absorber en un seul bloc.

La version 6 h se concentre sur \textbf{l'\'epine dorsale analytique} :

\begin{enumerate}
  \item \textbf{Python et Pandas} -- le socle pour les cinq modules suivants.
  \item \textbf{Nettoyage et statistique descriptive} -- la base de toute
        analyse cr\'edible.
  \item \textbf{Visualisation} -- communiquer un r\'esultat.
  \item \textbf{Tests d'hypoth\`eses} -- l'inf\'erence formelle.
  \item \textbf{R\'egression} -- l'ajustement multivariable.
  \item \textbf{Machine learning clinique} -- pr\'ediction, ROC, calibration,
        \'equit\'e.
\end{enumerate}

Trois sujets de la version 30 h ne sont \textbf{pas} couverts ici :
analyse temporelle approfondie, analyse spatiale (g\'eopandas), et projet
capstone. Une br\`eve mention de ces sujets en fin de document indique
o\`u les trouver.

\subsection*{Comment utiliser ce document avec les notebooks}

Le ratio recommand\'e est \textbf{15-20 minutes de lecture pour 40 minutes
de manipulation du notebook}. Pour chaque module :

\begin{enumerate}
  \item Lire la section correspondante de ces notes (environ 6-8 pages).
  \item Ouvrir le notebook \texttt{module\_0N\_*.ipynb} dans Jupyter ou
        Colab.
  \item Ex\'ecuter chaque cellule en lisant les commentaires.
  \item Faire les exercices, puis comparer avec les cellules
        \emph{Solution}.
\end{enumerate}

\begin{notebox}
Les notebooks sont rep\'er\'es par un en-t\^ete graphique qui indique le
num\'ero du module, la dur\'ee, et le niveau. La page d'accueil
\texttt{index.ipynb} liste les 6 modules avec un bouton
\og Ouvrir dans Colab \fg{} pour chacun.
\end{notebox}

\subsection*{Conventions typographiques}

\begin{itemize}
  \item \texttt{code} -- terme technique, nom de variable, fonction Python.
  \item \emph{italique} -- mise en relief d'un concept la premi\`ere fois
        qu'il est introduit.
  \item Les encarts color\'es signalent : \textbf{Astuce} (vert),
        \textbf{Attention} (orange), \textbf{Note} (bleu),
        \textbf{Contexte clinique} (vert plus dense).
\end{itemize}

\newpage

% ============================================================
% MODULE 1
% ============================================================

\section{Module 1 -- Python et Pandas pour la sant\'e}

\subsection{Pourquoi Python pour les donn\'ees de sant\'e}

Trois raisons font que Python s'est impos\'e comme l'outil par d\'efaut
dans la recherche en sant\'e contemporaine :

\textbf{Reproductibilit\'e.} Un script documente chaque \'etape de votre
analyse. Quand un examinateur demande \og comment avez-vous trait\'e les
valeurs de laboratoire manquantes ? \fg{}, vous montrez le code -- pas
une description de clics dans des menus. Les revues exp\'erimentales
s'orientent toutes dans cette direction.

\textbf{\'Echelle.} Excel commence \`a peiner sur 100\,000 lignes. Pandas
en traite des millions sans difficult\'e. Plus important encore :
\og l'analyse est s\'epar\'ee des donn\'ees \fg{}, vous pouvez ex\'ecuter
le m\^eme script sur un \'echantillon de 1 000 patients puis sur la base
enti\`ere sans rien changer.

\textbf{\'Ecosyst\`eme.} Des biblioth\`eques telles que \texttt{lifelines}
(analyse de survie), \texttt{statsmodels} (mod\`eles statistiques),
\texttt{scikit-learn} (machine learning), \texttt{geopandas}
(cartographie), \texttt{pingouin} (statistiques cliniques en fran\c{c}ais)
offrent des outils que SPSS ou Stata n'\'egalent pas. Et le tout est
gratuit.

\subsection{Variables et les quatre types fondamentaux}

Une \emph{variable} est un nom qui pointe vers une valeur. Vous n'avez
\textbf{pas \`a d\'eclarer} le type : Python le devine \`a partir de la
valeur affect\'ee.

\begin{minted}{python}
age = 45                           # int   -- nombre entier
temperature = 37.2                 # float -- nombre a virgule
diagnostic = "Diabete de type 2"   # str   -- texte entre guillemets
hospitalise = True                 # bool  -- True ou False
\end{minted}

Les quatre types couvrent l'essentiel des donn\'ees de sant\'e :

\begin{table}[h]
\centering
\begin{tabular}{lll}
\toprule
\bfseries Type & \bfseries Stocke & \bfseries Exemple clinique \\
\midrule
\texttt{int}   & Nombre entier         & \^Age, nombre d'admissions, parit\'e \\
\texttt{float} & Nombre \`a virgule    & Temp\'erature, HbA1c, cr\'eatinine \\
\texttt{str}   & Texte                 & Code CIM-10, identifiant patient \\
\texttt{bool}  & Vrai / Faux           & Fumeur ? Hospitalis\'e ? Grossesse ? \\
\bottomrule
\end{tabular}
\end{table}

\subsection{Collections : listes et dictionnaires}

Une \emph{liste} range plusieurs valeurs \textbf{dans l'ordre}. Une
s\'erie de tensions chez un m\^eme patient, par exemple, est naturellement
une liste.

\begin{minted}{python}
tensions = [120, 135, 128, 142, 118]   # mmHg, syst.
\end{minted}

L'indexation commence \`a \textbf{z\'ero} : \texttt{tensions[0]} est la
premi\`ere valeur, \texttt{tensions[-1]} la derni\`ere. La syntaxe
\texttt{tensions[1:3]} extrait les indices 1 et 2 (la borne droite est
exclue, c'est une convention math\'ematique cl\'e en Python).

Un \emph{dictionnaire} associe des \textbf{cl\'es} \`a des \textbf{valeurs}
-- exactement la forme d'une fiche patient.

\begin{minted}{python}
patient = {
    "id": "PAT-0042",
    "age": 58,
    "sexe": "F",
    "tension_systolique": 148,
    "medicaments": ["Amlodipine", "HCTZ"],
}
patient["age"]   # -> 58
\end{minted}

\begin{notebox}
Une \textbf{cohorte} = une \textbf{liste de dictionnaires}. C'est
exactement la forme que pandas convertit en \texttt{DataFrame} : la
fonction \texttt{pd.DataFrame(cohorte)} construit le tableau directement.
\end{notebox}

\subsection{Conditions et boucles : appliquer une r\`egle clinique}

Une \emph{conditionnelle} encode une r\`egle de d\'ecision. Une
\emph{boucle} l'applique \`a chaque patient d'une cohorte.

\begin{minted}{python}
for p in cohorte:
    if p["hba1c"] >= 6.5:
        statut = "Diabetique"
    elif p["hba1c"] >= 5.7:
        statut = "Pre-diabetique"
    else:
        statut = "Normal"
\end{minted}

L'\textbf{indentation} -- quatre espaces -- est la fa\c{c}on de Python de
d\'elimiter les blocs. Pas d'accolades, contrairement \`a C, Java ou R.
Une indentation incorrecte l\`eve une \texttt{IndentationError}.

L'ordre des conditions importe. On teste d'abord la plus stricte
($\geq 6{,}5$), puis la suivante. L'invers\'e classerait tout patient
avec HbA1c $\geq 5{,}7$ comme \og pr\'e-diab\'etique \fg{} -- y compris
les diab\'etiques av\'er\'es.

\subsection{Fonctions : emballer une formule clinique}

Une \emph{fonction} est un bloc de code nomm\'e, qui prend des
\textbf{arguments} et renvoie une \textbf{valeur}. Une fois \'ecrite,
on l'utilise comme une formule.

\begin{minted}{python}
def imc(poids_kg: float, taille_m: float) -> float:
    """Indice de Masse Corporelle en kg/m^2."""
    return poids_kg / (taille_m ** 2)
\end{minted}

Trois bonnes pratiques :

\begin{enumerate}
  \item \textbf{Docstring}. La cha\^ine entre triple guillemets juste
        apr\`es \texttt{def} documente la fonction. \texttt{help(imc)}
        l'affiche.
  \item \textbf{Annotations de type}
        (\texttt{poids\_kg: float}). Elles documentent l'intention
        sans rien forcer \`a l'ex\'ecution. Vivement recommand\'ees en
        contexte clinique.
  \item \textbf{Appels par mot-cl\'e}. Pr\'ef\'erez
        \texttt{imc(poids\_kg=82, taille\_m=1{.}75)} \`a
        \texttt{imc(82, 1{.}75)}.
\end{enumerate}

\begin{attention}
Pour toute fonction de calcul m\'edical (dosage, eGFR, scores), les
appels positionnels sont une source d'erreur grave. Une inversion entre
\texttt{poids\_kg} et \texttt{age} dans une formule de dosage mg/kg peut
\^etre dramatique. \textbf{Toujours} pr\'ef\'erer les appels par mot-cl\'e.
\end{attention}

\subsection{Pandas : passer du Python pur \`a la table}

\emph{Pandas} est la biblioth\`eque standard pour manipuler des tableaux
en Python. L'objet central est le \texttt{DataFrame} : pensez \`a un
tableur, mais programmable.

\begin{minted}{python}
import pandas as pd
gm = pd.read_csv("gapminder.tsv", sep="\t")
\end{minted}

Avant toute analyse, ex\'ecuter \textbf{les cinq m\'ethodes d'inspection}
fait gagner des heures :

\begin{table}[h]
\centering
\begin{tabular}{ll}
\toprule
\bfseries M\'ethode & \bfseries Ce qu'elle vous dit \\
\midrule
\texttt{df.head(n)}      & Les $n$ premi\`eres lignes \\
\texttt{df.shape}        & (nb lignes, nb colonnes) \\
\texttt{df.dtypes}       & Type de chaque colonne \\
\texttt{df.describe()}   & Moyenne, $\sigma$, quartiles \\
\texttt{df['x'].value\_counts()} & Fr\'equences d'une colonne cat\'egorielle \\
\bottomrule
\end{tabular}
\end{table}

\subsection{Filtrer, trier, grouper}

Le \textbf{filtrage} utilise un \emph{masque bool\'een} -- une colonne de
\texttt{True}/\texttt{False} qu'on place entre crochets.

\begin{minted}{python}
afrique = gm[gm["continent"] == "Africa"]
\end{minted}

Pour combiner deux conditions, utilisez \texttt{\&} (et), \texttt{|} (ou),
\texttt{\~{}} (non), et \textbf{paranth\'esez chaque condition}. Cette
derni\`ere obligation est non-n\'egociable : sans parenth\`eses,
Python applique les op\'erateurs dans le mauvais ordre.

\begin{minted}{python}
pauvre_africain = gm[
    (gm["continent"] == "Africa")
    & (gm["gdpPercap"] < 1000)
]
\end{minted}

Le \textbf{tri} se fait avec \texttt{.sort\_values("col")}. Pour obtenir
les $n$ plus grandes/petites valeurs directement,
\texttt{.nlargest(n, "col")} / \texttt{.nsmallest(n, "col")}.

Le \textbf{groupby} est le \emph{geste pivot} de pandas. Le patron est
\textbf{\og split, apply, combine \fg{}} (McKinney, 2017) :

\begin{enumerate}
  \item \textbf{Split} -- couper les lignes en groupes selon une
        cat\'egorie.
  \item \textbf{Apply} -- calculer une statistique sur chaque groupe.
  \item \textbf{Combine} -- recoller le tout en une table.
\end{enumerate}

\begin{minted}{python}
gm.groupby("continent")["lifeExp"].mean()
\end{minted}

\subsection{Cr\'eer une colonne calcul\'ee}

Pandas est \emph{vectoris\'e} : on ne fait pas de boucle pour remplir une
colonne. L'op\'eration est appliqu\'ee \`a toutes les lignes en m\^eme
temps.

\begin{minted}{python}
gm["gdp_milliards"] = gm["pop"] * gm["gdpPercap"] / 1e9
\end{minted}

Pour transformer une variable continue en cat\'egories ordonn\'ees,
\texttt{pd.cut} :

\begin{minted}{python}
gm["categorie_ev"] = pd.cut(
    gm["lifeExp"],
    bins=[0, 50, 65, 75, 100],
    labels=["Tres faible", "Faible", "Moyen", "Eleve"],
)
\end{minted}

\subsection{Points cl\'es du Module 1}

\begin{itemize}
  \item Quatre types Python suffisent pour 95 \% des donn\'ees cliniques.
  \item Une cohorte = liste de dictionnaires = \texttt{DataFrame}.
  \item Pandas se manipule par \emph{m\'ethodes encha\^in\'ees} :
        \texttt{df.filter(...).groupby(...).agg(...)}.
  \item Les masques bool\'eens demandent toujours des \textbf{parenth\`eses}
        autour de chaque condition.
  \item \texttt{pd.cut} est l'outil universel pour cat\'egoriser une
        variable continue (seuils cliniques, tranches d'\^age).
\end{itemize}

\newpage

% ============================================================
% MODULE 2
% ============================================================

\section{Module 2 -- Nettoyage et statistique descriptive}

\subsection{Le nettoyage des donn\'ees est un raisonnement clinique}

Aucune base de donn\'ees r\'eelle n'est propre. Dans les dossiers patient
informatis\'es (DPI), une \'etude de r\'ef\'erence
\citep{weiskopf2013methods} estime qu'environ 18 \% des valeurs de
laboratoire pr\'esentent un probl\`eme de qualit\'e : doublon,
incoh\'erence d'unit\'e, valeur impossible.

Nettoyer, c'est \textbf{prendre des d\'ecisions}. Supprimer une ligne ?
Imputer une valeur ? Marquer un \emph{flag} ? Chaque choix d\'eplace
l'analyse. La discipline consiste \`a \textbf{documenter chaque
d\'ecision} pour qu'un relecteur puisse les reproduire.

\subsection{Les trois m\'ecanismes de donn\'ees manquantes}

Le cadre conceptuel a \'et\'e \'etabli par \citet{rubin1976inference} et
n'a pas chang\'e depuis. Trois m\'ecanismes :

\textbf{MCAR -- Missing Completely At Random.} Le manquement est
ind\'ependant de tout. Un tube de pr\'el\`evement qui casse, par exemple.
Supprimer les lignes concern\'ees est \textbf{sans biais} ; on perd
seulement de la puissance.

\textbf{MAR -- Missing At Random.} Le manquement d\'epend de variables
\textbf{observ\'ees}. Exemple : le cholest\'erol n'est pas mesur\'e chez
les patients jeunes (l'\^age est observ\'e). L'imputation
\textbf{conditionnelle \`a l'\^age} est sans biais.

\textbf{MNAR -- Missing Not At Random.} Le manquement d\'epend de la
valeur \textbf{non observ\'ee} elle-m\^eme. Les patients les plus malades
ratent leur consultation de suivi, donc on n'a pas leur tension a
six mois. \textbf{C'est le cas le plus dangereux} : supprimer biaise
syst\'ematiquement vers les patients en meilleure sant\'e.

\begin{attention}
Dans les DPI, MNAR est la \textbf{r\`egle}, pas l'exception. Avant
d'imputer ou de supprimer, v\'erifiez \textbf{toujours} si l'issue (la
variable dont vous voulez parler) diff\`ere entre les lignes compl\`etes
et les lignes \`a donn\'ees manquantes. Si oui, le manquement est
informatif : il faut le mod\'eliser, pas le balayer.
\end{attention}

\subsection{D\'etecter et compter les valeurs manquantes}

Pandas utilise \texttt{NaN} (Not a Number) pour le num\'erique manquant
et \texttt{None} pour les objets. Les deux sont d\'etect\'es par
\texttt{.isnull()}.

\begin{minted}{python}
df.isnull().sum()          # nb manquant par colonne
df.isnull().any(axis=1)    # True/False par ligne
\end{minted}

Un point souvent n\'eglig\'e : v\'erifier les manquements
\textbf{simultan\'es}. La corr\'elation de la matrice
\texttt{df.isnull()} r\'ev\`ele si deux colonnes manquent
\textbf{ensemble}, indice possible d'une cause commune (la m\^eme
consultation rat\'ee, le m\^eme automate en panne).

\subsection{Strat\'egies d'imputation}

\begin{table}[h]
\centering
\small
\begin{tabular}{ll}
\toprule
\bfseries M\'ethode & \bfseries Quand l'utiliser \\
\midrule
M\'ediane           & Continu asym\'etrique (lab, tension, IMC) \\
Moyenne             & Continu sym\'etrique (taille adulte) \\
Mode                & Cat\'egoriel (sexe, statut tabagique) \\
Par sous-groupe     & Valeur attendue d\'epend d'un groupe connu \\
KNN                 & Plusieurs colonnes corr\'el\'ees manquent ensemble \\
MICE / IterativeImputer & Imputation multivari\'ee, plusieurs runs \\
\bottomrule
\end{tabular}
\end{table}

R\`egle pratique : \textbf{m\'ediane par d\'efaut} pour tout ce qui est
clinique. La majorit\'e des variables de laboratoire ont des queues
droites (cr\'eatinine, prot\'eine C-r\'eactive, ferritine). La moyenne y
est tir\'ee par les extr\^emes.

Pour les variables critiques (issue principale, exposition cl\'e), gardez
un \textbf{indicateur de manquement} \emph{avant} d'imputer :

\begin{minted}{python}
df["tension_manquait"] = df["tension"].isnull().astype(int)
df["tension"] = df["tension"].fillna(df["tension"].median())
\end{minted}

L'indicateur peut ensuite servir de pr\'edicteur dans un mod\`ele (le
fait que la valeur ait \'et\'e absente peut \^etre pronostique en soi).

\subsection{Validation par plages cliniquement plausibles}

Toute base de donn\'ees a des valeurs \textbf{absurdes} : un \^age de
$-3$ ans, une tension de 450 mmHg, une HbA1c de 55 (lue comme NGSP
alors qu'elle est en IFCC). Codez vos plages de plausibilit\'e en tant
que dictionnaire, appliquez-les syst\'ematiquement :

\begin{minted}{python}
PLAGES = {
    "age":     (0, 120),
    "tension": (60, 300),
    "hba1c":   (2.0, 20.0),
}
\end{minted}

Les valeurs hors plage sont \textbf{presque toujours} des erreurs de
saisie -- soit une faute de frappe (450 au lieu de 145), soit une
confusion d'unit\'e (5{,}5 mmol/L de glyc\'emie lu comme 5{,}5 mg/dL).

\begin{clinique}
Le cas des unit\'es m\'erite une mention sp\'eciale. La glyc\'emie en
mg/dL et mmol/L sont \`a un facteur 18 pr\`es : 100 mg/dL $\approx$ 5{,}5
mmol/L. Sans colonne d'unit\'e explicite, vous ne pouvez pas distinguer
\og 7 \fg{} pour une glyc\'emie a jeun
\og normale en mmol/L \fg{} de \og pathologique en mg/dL \fg{}. La
v\'erit\'e absolue : \textbf{exigez une colonne unit\'e d\`es la collecte}.
\end{clinique}

\subsection{Tendance centrale et dispersion}

\textbf{Moyenne ou m\'ediane ?} R\`egle pragmatique : si une valeur
extr\^eme d\'eplace la moyenne de plus de quelques pourcents, rapportez
la \textbf{m\'ediane}. La m\'ediane est robuste : une valeur aberrante \`a
340 mg/dL au milieu de glyc\'emies de 90-110 ne la d\'eplace presque
pas.

\textbf{Pour la dispersion}, l'\'ecart-type ($\sigma$) est l'outil par
d\'efaut pour les distributions sym\'etriques. Pour les distributions
\emph{asym\'etriques} (dur\'ee de s\'ejour, co\^uts hospitaliers, biologie
\`a longue queue), l'\textbf{intervalle interquartile (IQR = Q3 -- Q1)}
est plus honn\^ete.

Le \textbf{coefficient de variation} $CV = \sigma / \mu$ (en \%) est
sans unit\'e. Utile pour comparer la dispersion relative de variables
aux unit\'es diff\'erentes :
\og l'h\'emoglobine varie moins ($CV \approx 5\,\%$) que la
glyc\'emie ($CV \approx 8\,\%$) chez un m\^eme patient \fg{}.

\subsection{Mesures \'epid\'emiologiques de base}

Quatre mesures couvrent l'essentiel des rapports en sant\'e publique :

\textbf{Pr\'evalence} -- proportion de cas \textbf{existants} \`a un
instant $t$ :
\[
P = \frac{\text{nombre de cas}}{\text{population totale}}.
\]
Rapporter toujours avec un \textbf{intervalle de confiance}. Pour une
proportion proche de $0$ ou $1$, pr\'ef\'erer le \emph{Wilson} ou
\emph{Agresti-Coull} \`a l'approximation normale (Wald).

\textbf{Taux d'incidence} -- nouveaux cas par personne-temps :
\[
IR = \frac{\text{nouveaux cas}}{\text{personnes-ann\'ees suivies}}.
\]
Pour des \'ev\'enements rares, le compte suit une loi de Poisson, d'o\`u
$\mathrm{SE}(IR) = \sqrt{cas} / PA$.

\textbf{Ratio de risque (RR)} et \textbf{ratio de cotes (OR)} \`a partir
d'un tableau 2$\times$2 :

\begin{center}
\begin{tabular}{lcc}
\toprule
              & Malade$+$ & Malade$-$ \\
\midrule
Expos\'e      & $a$       & $b$ \\
Non expos\'e  & $c$       & $d$ \\
\bottomrule
\end{tabular}
\end{center}

\[
RR = \frac{a/(a+b)}{c/(c+d)}, \qquad
OR = \frac{a \cdot d}{b \cdot c}.
\]

L'OR est l'estim\'e dans les \'etudes \textbf{cas-t\'emoin} ; le RR dans
les \'etudes de \textbf{cohorte}. Pour de petits effets et issues rares,
$OR \approx RR$. D\`es que la pr\'evalence d\'epasse $\sim 10\,\%$, l'OR
\emph{surestime} le RR -- rapportez-le comme OR, pas comme RR.

Les intervalles de confiance pour RR et OR sont calcul\'es sur l'\'echelle
$\ln$ (la distribution y est approximativement normale), puis
exponentialis\'es.

\subsection{Pipeline de nettoyage r\'eutilisable}

Le bon r\'eflexe : bundler les \'etapes dans \textbf{une seule fonction},
sur le mod\`ele entr\'ee $\to$ sortie. Test\'ee une fois, elle s'applique
\`a tous les jeux de donn\'ees similaires.

\begin{minted}{python}
def nettoyer(df, plages=PLAGES):
    df = df.copy()
    df.columns = df.columns.str.strip().str.lower()
    df = df.drop_duplicates()
    for col, (lo, hi) in plages.items():
        if col in df.columns:
            df.loc[(df[col] < lo) | (df[col] > hi), col] = np.nan
    for col in df.select_dtypes(include="number").columns:
        df[col] = df[col].fillna(df[col].median())
    return df
\end{minted}

\subsection{Points cl\'es du Module 2}

\begin{itemize}
  \item Les trois m\'ecanismes MCAR / MAR / MNAR d\'eterminent ce qu'on
        peut faire sans biais.
  \item En DPI, MNAR est la r\`egle : v\'erifier que les lignes
        compl\`etes et incompl\`etes ont la m\^eme issue.
  \item M\'ediane par d\'efaut pour le clinique ; IQR pour la dispersion
        asym\'etrique.
  \item Pr\'evalence, incidence, RR, OR -- toujours rapporter avec IC.
  \item Encapsuler le nettoyage dans une fonction.
\end{itemize}

\newpage

% ============================================================
% MODULE 3
% ============================================================

\section{Module 3 -- Visualisation pour donn\'ees de sant\'e}

\subsection{Une figure r\'epond \`a une question}

Avant d'\'ecrire la moindre ligne de code de visualisation, formulez la
\textbf{question} \`a laquelle la figure r\'epond. \og Distribution de
l'IMC \fg{} n'est pas une question. \og La distribution de l'IMC
diff\`ere-t-elle entre les groupes d'\^age ? \fg{} en est une.

Le type de question d\'etermine le type de graphique.

\begin{table}[h]
\centering
\small
\begin{tabular}{ll}
\toprule
\bfseries Question & \bfseries Type de graphique \\
\midrule
Comment une variable est-elle distribu\'ee ?      & Histogramme, KDE \\
Comment une variable diff\`ere-t-elle entre groupes ? & Boite \`a moustaches \\
Comment compare-t-on des comptages ?              & Barres \\
Comment \'evolue-t-elle dans le temps ?           & Ligne \\
Deux variables continues sont-elles associ\'ees ? & Nuage de points \\
Quelles variables sont corr\'el\'ees ?            & Heatmap \\
\bottomrule
\end{tabular}
\end{table}

\subsection{R\`egles non-n\'egociables}

Toute figure publi\'ee ou pr\'esent\'ee en r\'eunion respecte ces r\`egles
de base. Les violer trahit un manque de rigueur que vos lecteurs
remarqueront :

\begin{enumerate}
  \item \textbf{Les axes portent leur unit\'e.} \og Tension systolique
        (mmHg) \fg{}, jamais juste \og BP \fg{}.
  \item \textbf{Les barres commencent \`a z\'ero.} Tronquer l'axe
        vertical pour amplifier visuellement une diff\'erence est l'un
        des mensonges visuels les plus document\'es \citep{tufte2001}.
  \item \textbf{Pas de 3D.} Aucun graphique 3D n'a jamais am\'elior\'e la
        lisibilit\'e. Ils introduisent des distorsions de perception.
  \item \textbf{Pas de camemberts au-del\`a de 3 parts.} L'\oe il humain
        compare mal les angles ; les barres font mieux.
  \item \textbf{Pas de double axe Y.} Sauf cas exceptionnel, ils
        sugg\`erent fausse corr\'elation et permettent de manipuler la
        perception.
  \item \textbf{Palette accessible aux daltoniens.} Suivez les
        recommandations de \citet{wong2011points}.
  \item \textbf{Annoter} la taille d'\'echantillon ($n = 1\,234$) et les
        \textbf{seuils cliniques} (lignes horizontales / verticales aux
        cut-offs pertinents).
\end{enumerate}

\subsection{Histogrammes et estimation de densit\'e}

L'histogramme montre comment une variable continue est r\'epartie. Le
nombre de classes (\texttt{bins}) compte : trop peu cache la structure,
trop transforme le graphe en bruit. Une r\`egle empirique :
$k \approx \sqrt{n}$.

\begin{minted}{python}
ax.hist(donnees, bins=20, edgecolor="white")
ax.axvline(126, color="red", ls="--",
           label="Seuil diabete (126 mg/dL)")
\end{minted}

Pour comparer plusieurs distributions sur les m\^emes axes, l'estimation
\emph{kernel density} (KDE) est plus lisible : courbes liss\'ees au lieu
de barres se chevauchant.

\subsection{Boites \`a moustaches}

Anatomie : la bo\^ite va de $Q_1$ \`a $Q_3$, la barre m\'ediane est le
$Q_2$, les moustaches s'\'etendent au dernier point dans
$1{,}5 \times IQR$, les points au-del\`a sont les valeurs aberrantes
individuelles.

Limite importante : \textbf{les bo\^ites cachent la bimodalit\'e}. Deux
distributions tr\`es diff\'erentes peuvent produire la m\^eme bo\^ite. En
cas de doute, superposer un \emph{stripplot} ou montrer l'histogramme en
parall\`ele.

\subsection{Lignes : tendances temporelles et courbes \'epid\'emiques}

Le patron canonique pour une \'epid\'emie : \textbf{barres journali\`eres}
+ \textbf{moyenne mobile sur 7 jours} en haut, \textbf{courbe cumul\'ee}
en bas. La barre montre le bruit r\'eel, la ligne le signal liss\'e.

La moyenne mobile sur 7 jours efface la p\'eriodicit\'e hebdomadaire
(d\'eclarations rep\'erables en bloc le lundi, vacances scolaires, etc.).

\begin{minted}{python}
epi["mm_7j"] = epi["nouveaux_cas"].rolling(7).mean()
\end{minted}

\subsection{Nuages de points et associations}

Le scatter plot affiche deux variables continues sur un plan. Pour
encoder une troisi\`eme variable, deux choix s'offrent : la \textbf{taille
des points} (bubble chart) ou la \textbf{couleur}. La taille marche bien
pour des grandeurs sur plusieurs ordres -- la population des pays, par
exemple.

L'\'echelle \textbf{logarithmique} sur un axe est cruciale d\`es que la
variable s'\'etale sur 3-4 ordres de grandeur (PIB par habitant, charges
virales, comptages bact\'eriens). Sans le log, 95 \% des points
s'agglutinent dans un coin.

\subsection{Le tableau de bord multi-panneaux}

Un tableau de bord est une \textbf{seule figure} avec plusieurs
panneaux qui r\'epondent ensemble \`a une question g\'en\'erale (\og que
sait-on de la sant\'e mondiale en 2007 ? \fg{}).

R\`egles compl\'ementaires pour les multi-panneaux :

\begin{enumerate}
  \item \textbf{Lettrer} les panneaux (A, B, C, D) pour les r\'ef\'erencer
        dans la l\'egende et le texte.
  \item \textbf{Garder le m\^eme code couleur} entre panneaux : un
        continent gard\'e en bleu dans le panneau A doit rester bleu en B.
  \item \textbf{Hi\'erarchie de l'\oe il} : panneau A en haut \`a gauche
        re\c{c}oit l'\oe il en premier ; mettez-y l'image la plus
        importante.
\end{enumerate}

\subsection{Points cl\'es du Module 3}

\begin{itemize}
  \item Le type de graphique d\'epend du type de question, pas de
        l'\'esth\'etique.
  \item Palette colour-blind-safe par d\'efaut (Wong 2011).
  \item Annoter unit\'es, $n$, seuils cliniques.
  \item Bo\^ite \`a moustaches : utile pour comparer, dangereuse pour
        l'unimodalit\'e.
  \item Tableau de bord = une question, plusieurs panneaux compl\'ementaires.
\end{itemize}

\newpage

% ============================================================
% MODULE 4
% ============================================================

\section{Module 4 -- Tests d'hypoth\`eses}

\subsection{La logique d'un test}

Un test d'hypoth\`ese est une \emph{preuve par contradiction
probabiliste} :

\begin{enumerate}
  \item On suppose la \textbf{nulle $H_0$ vraie} (pas d'effet).
  \item On calcule une statistique de test \`a partir des donn\'ees.
  \item On se demande : \emph{si $H_0$ \'etait vraie, ces donn\'ees
        seraient-elles inhabituelles ?} Cette probabilit\'e est la
        \textbf{p-value}.
  \item Si la p-value est $< \alpha$ (souvent 0,05), on \textbf{rejette
        $H_0$}.
\end{enumerate}

\subsection{Ce qu'une p-value n'est pas}

Aucune confusion n'a co\^ut\'e plus cher \`a la recherche en sant\'e que
les interpr\'etations erron\'ees de la p-value. La d\'eclaration officielle
de l'American Statistical Association \citep{wasserstein2016asa} en
\'enum\`ere les plus dommageables. R\'ep\'etons les essentielles :

\begin{itemize}
  \item Ce n'est \textbf{pas} la probabilit\'e que $H_0$ soit vraie.
  \item Ce n'est \textbf{pas} la probabilit\'e que le r\'esultat soit
        d\^u au hasard.
  \item Elle ne mesure \textbf{pas} la taille d'un effet.
  \item Un test significatif sur un grand $n$ peut correspondre \`a un
        effet \textbf{cliniquement n\'egligeable}.
\end{itemize}

\begin{attention}
La r\`egle d'or : \textbf{toujours rapporter une p-value avec une taille
d'effet et un intervalle de confiance}. Sur 100\,000 patients, un essai
peut d\'etecter une diff\'erence moyenne de 0,5 mmHg avec
$p < 10^{-6}$. Statistiquement significatif, cliniquement non pertinent.
\end{attention}

\subsection{Test t : un, deux, ou pairs}

\textbf{Test t \`a un \'echantillon}. La moyenne d'un \'echantillon
diff\`ere-t-elle d'une valeur de r\'ef\'erence connue ?
Exemple : le cholest\'erol moyen dans notre clinique diff\`ere-t-il de la
moyenne nationale de 200 mg/dL ?

\textbf{Test t \`a deux \'echantillons ind\'ependants}. Deux groupes
ind\'ependants diff\`erent-ils ? Exemple : la tension dans le groupe
trait\'e diff\`ere-t-elle de celle du placebo ?

Conditions : ind\'ependance, normalit\'e dans chaque groupe (assoupli par
le TCL pour $n \gtrsim 30$), \textbf{\'egalit\'e des variances}. Cette
derni\`ere se rel\^ache avec la \textbf{correction de Welch}
(\texttt{equal\_var=False}), qui est le \textbf{d\'efaut moderne
recommand\'e}. Pratique : toujours mettre \texttt{equal\_var=False}
\citep{ruxton2006unequal}.

\textbf{Test t appari\'e}. Les m\^emes patients mesur\'es avant et apr\`es
une intervention. Le test utilise les \textbf{diff\'erences
intra-patient}, \'eliminant la variabilit\'e inter-patient et augmentant la
puissance.

\begin{minted}{python}
from scipy import stats
t, p = stats.ttest_ind(traitement, controle, equal_var=False)
\end{minted}

\subsection{Test du chi-deux d'ind\'ependance}

Pour deux variables \textbf{cat\'egorielles}, le chi-deux compare les
comptages observ\'es aux comptages attendus sous ind\'ependance :
\[
\chi^2 = \sum_{ij} \frac{(O_{ij} - E_{ij})^2}{E_{ij}}.
\]

Condition : chaque effectif attendu doit \^etre $\geq 5$. Sinon utiliser
le \textbf{test exact de Fisher} (\texttt{stats.fisher\_exact}).

\subsection{L'alternative non param\'etrique : Mann-Whitney U}

Quand la variable est asym\'etrique (dur\'ee de s\'ejour, co\^uts) ou que
$n$ est petit, le test t sur l'\'echelle brute peut induire en erreur.
Mann-Whitney U compare les \textbf{rangs} : il teste si une distribution
domine stochastiquement l'autre.

\begin{notebox}
Comparer Mann-Whitney U et un test t apr\`es transformation log :
souvent les deux donnent des conclusions similaires sur les donn\'ees
asym\'etriques. La transformation log a l'avantage de rendre les
coefficients de r\'egression interpr\'etables (en pourcentage de
changement).
\end{notebox}

\subsection{Correction pour tests multiples}

Si vous testez \textbf{20 hypoth\`eses ind\'ependantes au seuil 0{,}05
sous $H_0$}, vous attendez en moyenne \textbf{un faux positif par
hasard}. Sur 1\,000 tests sans effet r\'eel, ~50 sont \og significatifs \fg{}
\`a $p < 0{,}05$. C'est l'\emph{erreur familiale}.

Deux familles de corrections :

\textbf{Bonferroni} (contr\^ole le \emph{family-wise error rate}) :
multiplier chaque p par $m$ (ou comparer \`a $\alpha / m$). Tr\`es
conservateur ; appropri\'e quand m\^eme \textbf{un} faux positif est
inacceptable (essais cliniques, d\'ecisions r\'eglementaires).

\textbf{Benjamini-Hochberg FDR} \citep{benjamini1995controlling} :
contr\^ole la proportion attendue de faux positifs parmi les positifs
d\'eclar\'es. Standard en g\'enomique et en analyses exploratoires.

\begin{minted}{python}
from statsmodels.stats.multitest import multipletests
_, p_corr, _, _ = multipletests(pvals, method="fdr_bh")
\end{minted}

\subsection{Taille d'effet}

Le d de Cohen :
\[
d = \frac{\bar{x}_1 - \bar{x}_2}{s_{pooled}}.
\]

Rep\`eres indicatifs \citep{cohen1988statistical} :

\begin{center}
\small
\begin{tabular}{ll}
\toprule
$|d|$ & Interpr\'etation \\
\midrule
$< 0{,}2$  & N\'egligeable \\
$0{,}2$ - $0{,}5$ & Petit \\
$0{,}5$ - $0{,}8$ & Moyen \\
$\geq 0{,}8$ & Large \\
\bottomrule
\end{tabular}
\end{center}

Ces rep\`eres ont \'et\'e \'etablis dans le domaine des sciences sociales
et doivent \^etre interpr\'et\'es avec prudence en clinique : une
diff\'erence \og petite \fg{} sur la mortalit\'e cardiovasculaire peut
\^etre tr\`es importante au niveau populationnel.

\subsection{Points cl\'es du Module 4}

\begin{itemize}
  \item p-value $\neq$ probabilit\'e que $H_0$ soit vraie.
  \item Toujours rapporter taille d'effet \emph{et} IC \`a c\^ot\'e de la
        p-value.
  \item Welch (\texttt{equal\_var=False}) par d\'efaut pour le test t \`a
        deux \'echantillons.
  \item Mann-Whitney U pour les distributions asym\'etriques.
  \item Tests multiples : Bonferroni pour ne tol\'erer aucun FP ;
        BH-FDR pour les screenings.
\end{itemize}

\newpage

% ============================================================
% MODULE 5
% ============================================================

\section{Module 5 -- R\'egression}

\subsection{De la corr\'elation \`a la pr\'ediction}

La r\'egression est l'outil le plus utilis\'e en recherche clinique
contemporaine. Elle r\'epond \`a \textbf{deux questions distinctes} qu'il
faut savoir s\'eparer.

\textbf{Inf\'erentielle} : quels facteurs sont associ\'es \`a une issue,
en ajustant pour les autres ? Exemple : l'IMC est-il associ\'e \`a la
tension, ind\'ependamment de l'\^age et du sexe ? L'outil de
pr\'edilection est \texttt{statsmodels}, qui donne coefficients,
IC 95\,\%, p-values, $R^2$.

\textbf{Pr\'edictive} : pour un nouveau patient, quelle est la
probabilit\'e d'issue ? Quel est son risque \`a 10 ans ? L'outil de
pr\'edilection est \texttt{scikit-learn}, qui donne pr\'edictions et
m\'etriques de performance.

\begin{notebox}
\texttt{statsmodels} et \texttt{scikit-learn} ne sont \textbf{pas
redondants}. Le m\^eme mod\`ele (r\'egression lin\'eaire ou logistique) ne
vous renverra pas la m\^eme chose : \texttt{statsmodels} pour
interpr\'eter, \texttt{scikit-learn} pour pr\'edire. Apprenez \`a
utiliser les deux.
\end{notebox}

\subsection{R\'egression lin\'eaire simple et multiple}

Le mod\`ele lin\'eaire :
\[
y_i = \beta_0 + \beta_1 x_{1i} + \beta_2 x_{2i} + \cdots + \varepsilon_i,
\quad
\varepsilon_i \sim \mathcal{N}(0, \sigma^2).
\]

Chaque coefficient $\beta_k$ s'interpr\`ete comme le \textbf{changement
attendu de $y$ pour une augmentation de une unit\'e de $x_k$, en
maintenant les autres pr\'edicteurs constants}. C'est la cl\'e de
l'ajustement pour facteur de confusion.

\subsection{Facteur de confusion : brute vs ajust\'ee}

Un facteur de confusion est associ\'e \`a l'\textbf{exposition} et \`a
l'\textbf{issue}, sans \^etre sur le chemin causal entre les deux. Si on
l'ignore, l'estimation brute est biais\'ee.

Exemple-\'ecole : le caf\'e semble \^etre un facteur de risque
cardiovasculaire dans une analyse brute. Mais une fois qu'on ajuste pour
le tabagisme, le coefficient du caf\'e s'\'evanouit. Le tabagisme \'etait
le \emph{vrai} facteur, et il \'etait associ\'e \`a la consommation de
caf\'e -- d'o\`u l'illusion.

La r\'egression multiple ajuste \emph{simultan\'ement} pour tous les
confondants observ\'es. Elle ne corrige \textbf{pas} les confondants
inobservs\'es : pour cela il faut des m\'ethodes causales avanc\'ees
(variables instrumentales, score de propension, \emph{difference-in-
differences}).

\subsection{V\'erifier les hypoth\`eses}

Quatre conditions pour des inf\'erences valides :

\begin{enumerate}
  \item \textbf{Lin\'earit\'e} -- la relation entre $x_k$ et $y$ est
        lin\'eaire. Diagnostic : graphique r\'esidus vs pr\'edits.
  \item \textbf{Ind\'ependance} -- les observations sont ind\'ependantes
        (faux si donn\'ees longitudinales : utiliser un mod\`ele mixte).
  \item \textbf{Homosc\'edasticit\'e} -- la variance des r\'esidus est
        constante. Un nuage en entonnoir est le drapeau rouge.
  \item \textbf{Normalit\'e des r\'esidus} -- diagnostic QQ-plot. Pas
        besoin d'\^etre parfait pour $n$ grand (TCL).
\end{enumerate}

\subsection{R\'egression logistique pour les issues binaires}

Quand l'issue est oui/non, l'OLS pr\'edit des valeurs en dehors de
$[0, 1]$ -- non utilisable. La logistique mod\'elise la
\textbf{probabilit\'e} :
\[
\Pr(y = 1 \mid x) = \sigma(\beta_0 + \beta_1 x_1 + \cdots),
\quad
\sigma(z) = \frac{1}{1 + e^{-z}}.
\]

Les coefficients sont des \textbf{log-odds}. Pour les transformer en
\textbf{odds ratios} interpr\'etables : $OR = e^{\beta}$. Un OR de 1{,}04
par mg/dL de glyc\'emie signifie \og chaque mg/dL multiplie les cotes de
diab\`ete par 1{,}04 \fg{}. Sur une plage de 30 mg/dL, l'effet compos\'e
est $1{,}04^{30} \approx 3$ -- substantiel.

\subsection{\'Evaluation : matrice de confusion et les quatre taux}

Sur un \'echantillon test, comparer pr\'edictions et v\'erit\'e dans une
matrice $2 \times 2$ :

\begin{center}
\begin{tabular}{lcc}
\toprule
                       & Pr\'edit + & Pr\'edit -- \\
\midrule
R\'eel +               & VP         & FN \\
R\'eel --              & FP         & VN \\
\bottomrule
\end{tabular}
\end{center}

Les quatre m\'etriques cl\'es :

\begin{itemize}
  \item \textbf{Sensibilit\'e} = $VP / (VP+FN)$. Capacit\'e \`a
        \textbf{d\'etecter} les malades. Cruciale pour les outils de
        \emph{rule-out} (\og si test n\'egatif, je sais que ce n'est pas
        \c{c}a \fg{}).
  \item \textbf{Sp\'ecificit\'e} = $VN / (VN+FP)$. Capacit\'e \`a
        \textbf{exclure} les non-malades. Cruciale pour les outils de
        \emph{rule-in}.
  \item \textbf{VPP} (valeur pr\'edictive positive) = $VP / (VP+FP)$.
        Quand le test est positif, quelle est la probabilit\'e d'\^etre
        vraiment malade.
  \item \textbf{VPN} = $VN / (VN+FN)$. Quand le test est n\'egatif,
        quelle est la probabilit\'e d'\^etre vraiment sain.
\end{itemize}

\begin{clinique}
Sensibilit\'e et sp\'ecificit\'e d\'ependent uniquement du \textbf{test}.
VPP et VPN d\'ependent aussi de la \textbf{pr\'evalence dans la
population test\'ee}. Le m\^eme test, appliqu\'e en population
g\'en\'erale (faible pr\'evalence) puis aux urgences (forte
pr\'evalence), donne des VPP / VPN tr\`es diff\'erents. C'est pourquoi
les recommandations de d\'epistage d\'ependent du contexte
\'epid\'emiologique local.
\end{clinique}

\subsection{Points cl\'es du Module 5}

\begin{itemize}
  \item \texttt{statsmodels} pour l'inf\'erence ; \texttt{scikit-learn}
        pour la pr\'ediction.
  \item Coefficient = effet net \og toutes choses \'egales par
        ailleurs \fg{}, parmi les variables observ\'ees.
  \item OR = $e^{\beta}$ pour la logistique.
  \item Toujours ajuster pour les confondants \textbf{observ\'es}.
  \item Sensibilit\'e / sp\'ecificit\'e d\'ependent du test ; VPP / VPN
        d\'ependent de la pr\'evalence.
\end{itemize}

\newpage

% ============================================================
% MODULE 6
% ============================================================

\section{Module 6 -- Machine learning pour la pr\'ediction clinique}

\subsection{Quand ne pas utiliser le ML}

Le machine learning est s\'eduisant. Il est aussi \textbf{surutilis\'e}.
Trois situations o\`u il est inappropri\'e :

\textbf{Petits jeux de donn\'ees ($n < 200$).} Les arbres et les for\^ets
surajustent. La r\'egression logistique avec quelques pr\'edicteurs
choisis cliniquement est plus fiable et g\'en\'eralise mieux.

\textbf{Validation r\'eglementaire.} La FDA, l'EMA, l'AFSSAPS et la
plupart des agences exigent des mod\`eles \textbf{interpr\'etables}. Un
mod\`ele bo\^ite-noire ne passera pas l'\'evaluation. La r\'egression
logistique avec quelques pr\'edicteurs bien justifi\'es est le standard
de fait.

\textbf{Une r\`egle simple marche d\'ej\`a.} \og Glyc\'emie $> 126$ mg/dL
\fg{} identifie l'essentiel des diab\'etiques. \og INR $> 4$ \fg{}
identifie le surdosage en AVK. Si une r\`egle a une bonne sensibilit\'e,
elle est plus simple, plus auditable et plus enseignable qu'un random
forest.

\begin{notebox}
Une revue syst\'ematique r\'ecente \citep{christodoulou2019systematic} a
montr\'e que sur 71 \'etudes comparant ML \`a la r\'egression logistique
en pr\'ediction clinique, \textbf{aucun b\'en\'efice consistent en
performance n'a \'et\'e d\'emontr\'e} pour le ML. La r\'egression
logistique reste un baseline tr\`es difficile \`a battre.
\end{notebox}

\subsection{Train / test stratifi\'e}

Toujours \textbf{stratifier} sur l'issue. Sans stratification, le jeu
test peut avoir une pr\'evalence tr\`es diff\'erente du jeu
d'entra\^inement et les m\'etriques deviennent instables.

\begin{minted}{python}
from sklearn.model_selection import train_test_split
X_tr, X_te, y_tr, y_te = train_test_split(
    X, y, test_size=0.3,
    random_state=42, stratify=y
)
\end{minted}

\subsection{Validation crois\'ee $k$-folds}

Pour les petits jeux, le test unique est bruyant. La \textbf{validation
crois\'ee} divise les donn\'ees en $k$ plis, entra\^ine sur $k-1$, teste
sur le restant, et moyenne les $k$ scores obtenus.

Cinq plis est le d\'efaut classique. Dix plis donne des estimations
moins biais\'ees mais plus bruyantes. Pour les outcomes rares (incidence
$< 5\,\%$), utiliser une \textbf{validation crois\'ee stratifi\'ee}
(\texttt{StratifiedKFold}) qui pr\'eserve la pr\'evalence dans chaque
pli.

\subsection{Arbres et for\^ets al\'eatoires}

Un \emph{arbre de d\'ecision} est une s\'equence de questions oui/non
aboutissant \`a une feuille avec une classe pr\'edite. Sa force : la
\textbf{lisibilit\'e clinique}. Un clinicien peut suivre l'arbre et
expliquer la pr\'ediction au lit du patient.

Sa faiblesse : sans contrainte de profondeur, l'arbre m\'emorise les
donn\'ees d'entra\^inement et g\'en\'eralise mal. \textbf{Toujours fixer}
\texttt{max\_depth}.

Une \emph{for\^et al\'eatoire} combine des centaines d'arbres
entra\^in\'es sur des \'echantillons \emph{bootstrap} avec des
sous-ensembles al\'eatoires de variables, puis moyenne leurs
pr\'edictions. La randomisation d\'ecorr\`ele les arbres et r\'eduit la
variance.

Trade-off : meilleures pr\'edictions, mais on perd la lisibilit\'e
d'\og un arbre \fg{} pointable.

\subsection{Courbe ROC et AUC}

La courbe ROC trace la \textbf{sensibilit\'e} contre $1 -
\textbf{sp\'ecificit\'e}$ quand on fait varier le seuil de d\'ecision de
0 \`a 1. L'\textbf{AUC} (Area Under Curve) r\'esume tout en un nombre.

\begin{table}[h]
\centering
\begin{tabular}{ll}
\toprule
AUC & Interpr\'etation \\
\midrule
$0{,}5$              & Pas mieux qu'un tirage \\
$0{,}7 - 0{,}8$      & Acceptable \\
$0{,}8 - 0{,}9$      & Excellent \\
$> 0{,}9$            & Exceptionnel (rare en clinique r\'eelle) \\
\bottomrule
\end{tabular}
\end{table}

\begin{attention}
L'AUC seule ne suffit pas. Un mod\`ele peut avoir une excellente AUC
\textbf{et \^etre inutile} si ses probabilit\'es sont mal calibr\'ees,
ou s'il pr\'esente un biais entre sous-groupes. \textbf{Toujours}
compl\'eter l'AUC par une courbe de calibration et un audit d'\'equit\'e.
\end{attention}

\subsection{Calibration : les probabilit\'es sont-elles dignes de confiance ?}

Une bonne AUC n'implique \textbf{pas} de bonnes probabilit\'es. Un
mod\`ele \textbf{bien calibr\'e} produit des probabilit\'es qui
correspondent aux fr\'equences observ\'ees : si vous pr\'edisez 30 \% de
risque \`a 10 ans, environ 30 \% des patients d\'eveloppent
effectivement la maladie.

On v\'erifie avec :

\begin{itemize}
  \item \textbf{Courbe de calibration} -- probabilit\'e moyenne pr\'edite
        vs proportion observ\'ee, par d\'ecile de risque. Une bonne
        calibration s'approche de la diagonale.
  \item \textbf{Score de Brier} -- erreur quadratique moyenne entre
        probabilit\'es pr\'edites et issues 0/1. Plus bas = mieux.
\end{itemize}

\subsection{Audit d'\'equit\'e}

V\'erifier que le mod\`ele performe \textbf{de mani\`ere comparable}
entre groupes (sexe, \^age, ethnie, comorbidit\'e). Un \'ecart de
plusieurs points d'AUC entre groupes est un signal d'alerte.

\citet{obermeyer2019dissecting} ont d\'emontr\'e qu'un algorithme
largement d\'eploy\'e aux \'Etats-Unis sous-allouait syst\'ematiquement
les soins aux patients afro-am\'ericains -- non par dessein, mais parce
que la d\'epense pass\'ee servait de proxy au besoin futur, et les
patients afro-am\'ericains d\'epensaient moins \`a niveau de morbidit\'e
\'egal. L'audit d'\'equit\'e est un \textbf{contr\^ole minimal} avant tout
d\'eploiement.

\subsection{Crit\`eres pour d\'eployer un mod\`ele}

Avant de mettre un mod\`ele en production clinique :

\begin{enumerate}
  \item \textbf{Validation externe} sur un jeu ind\'ependant (autre
        h\^opital, autre p\'eriode). La performance baisse souvent de
        plusieurs points d'AUC -- attendez-vous-y.
  \item \textbf{Audit d'\'equit\'e} d\'ej\`a discut\'e.
  \item \textbf{Calibration} explicite, avec recalibration si
        n\'ecessaire.
  \item \textbf{D\'efinition du seuil de d\'ecision} co-construite avec
        les cliniciens et li\'ee \`a une action concr\`ete.
  \item \textbf{Monitoring} en production pour d\'etecter le drift.
  \item \textbf{Documentation TRIPOD-AI} \citep{collins2024tripodai}.
\end{enumerate}

\subsection{Points cl\'es du Module 6}

\begin{itemize}
  \item Le ML n'est pas toujours le bon outil ; la logistique reste un
        baseline difficile \`a battre.
  \item Stratifier le train/test ; valider en CV $k$-folds.
  \item AUC est n\'ecessaire mais pas suffisant : ajouter calibration
        et audit d'\'equit\'e.
  \item Validation externe avant tout d\'eploiement clinique.
  \item Documenter selon TRIPOD-AI.
\end{itemize}

\newpage

% ============================================================
% POUR ALLER PLUS LOIN
% ============================================================

\section{Pour aller plus loin}

Cette version 6 heures se concentre sur l'\'epine dorsale analytique.
La version 30 heures couvre trois sujets non abord\'es ici :

\textbf{Analyse temporelle approfondie.} D\'ecomposition saisonni\`ere
des s\'eries de cas, mod\'elisation ARIMA, lissage exponentiel, courbes
\'epid\'emiques avec retards de d\'eclaration. Voir
\texttt{code/notebooks/ch09\_geospatial\_temporel.ipynb}.

\textbf{Analyse spatiale.} GeoPandas, projections, choropleths,
auto-corr\'elation spatiale, hotspot analysis avec les indices de Moran
et Getis-Ord. Voir le m\^eme chapitre.

\textbf{Projet capstone.} Analyse compl\`ete end-to-end sur un jeu r\'eel,
r\'edaction d'un rapport et soutenance orale. Voir
\texttt{code/notebooks/ch10\_projet\_final.ipynb}.

Ressources additionnelles incontournables :

\begin{itemize}
  \item \textbf{An Introduction to Statistical Learning} (James, Witten,
        Hastie, Tibshirani) -- gratuit en ligne :
        \url{https://www.statlearning.com/}.
  \item \textbf{Python Data Science Handbook} (Jake VanderPlas) --
        gratuit en ligne :
        \url{https://jakevdp.github.io/PythonDataScienceHandbook/}.
  \item \textbf{Regression Modeling Strategies} (Frank Harrell) -- la
        bible de la r\'egression en biostatistique.
  \item \textbf{TRIPOD-AI Statement} pour le rapportage de mod\`eles de
        pr\'ediction en sant\'e : \url{https://www.tripod-statement.org/}.
\end{itemize}

\newpage

% ============================================================
% ANNEXE : AIDE-MEMOIRE
% ============================================================

\section*{Annexe -- Aide-m\'emoire des snippets cl\'es}
\addcontentsline{toc}{section}{Annexe -- Aide-m\'emoire}

\subsection*{Charger un fichier}
\begin{minted}{python}
df = pd.read_csv("fichier.csv")
df = pd.read_csv("fichier.tsv", sep="\t")
df = pd.read_excel("fichier.xlsx", sheet_name=0)
\end{minted}

\subsection*{Inspecter}
\begin{minted}{python}
df.head(); df.shape; df.dtypes; df.describe()
df["col"].value_counts(); df["col"].nunique()
df.isnull().sum()
\end{minted}

\subsection*{Filtrer / trier / grouper}
\begin{minted}{python}
df[df["col"] > 5]
df[(df["a"] > 1) & (df["b"] == "X")]
df.sort_values("col", ascending=False)
df.nlargest(10, "col")
df.groupby("cat")["x"].agg(["mean", "median"])
\end{minted}

\subsection*{Cr\'eer une colonne}
\begin{minted}{python}
df["nouvelle"] = df["a"] * df["b"]
df["categ"] = pd.cut(df["x"], bins=[0, 1, 2], labels=["L", "H"])
\end{minted}

\subsection*{Nettoyer}
\begin{minted}{python}
df = df.drop_duplicates()
df = df.dropna(subset=["col"])
df["col"] = df["col"].fillna(df["col"].median())
df.loc[df["col"] < 0, "col"] = np.nan
\end{minted}

\subsection*{Visualiser (matplotlib + seaborn)}
\begin{minted}{python}
fig, ax = plt.subplots(figsize=(8, 5))
ax.hist(s, bins=30); ax.set_xlabel("..."); ax.set_ylabel("...")
sns.boxplot(data=df, x="cat", y="val", ax=ax)
ax.axvline(seuil, color="red", ls="--")
plt.tight_layout(); plt.show()
\end{minted}

\subsection*{Tests}
\begin{minted}{python}
from scipy import stats
stats.ttest_ind(a, b, equal_var=False)        # Welch
stats.ttest_rel(avant, apres)                  # apparie
stats.chi2_contingency(table)                  # chi2
stats.mannwhitneyu(g1, g2, alternative="two-sided")
\end{minted}

\subsection*{R\'egression}
\begin{minted}{python}
import statsmodels.api as sm
X = sm.add_constant(df[predicteurs])
ols = sm.OLS(df["y"], X).fit()
logit = sm.Logit(df["y_bin"], X).fit(disp=False)
np.exp(logit.params)                           # OR
\end{minted}

\subsection*{Machine learning}
\begin{minted}{python}
from sklearn.model_selection import train_test_split, cross_val_score
from sklearn.linear_model import LogisticRegression
from sklearn.ensemble import RandomForestClassifier
from sklearn.metrics import roc_auc_score, confusion_matrix

X_tr, X_te, y_tr, y_te = train_test_split(
    X, y, test_size=0.3, random_state=42, stratify=y)
rf = RandomForestClassifier(n_estimators=200, max_depth=6,
                            random_state=42).fit(X_tr, y_tr)
roc_auc_score(y_te, rf.predict_proba(X_te)[:, 1])
\end{minted}

\newpage

% ============================================================
% BIBLIOGRAPHIE
% ============================================================

\section*{Bibliographie}
\addcontentsline{toc}{section}{Bibliographie}

\renewcommand{\refname}{}

\begin{thebibliography}{99}

\bibitem[Benjamini \& Hochberg, 1995]{benjamini1995controlling}
Benjamini, Y. \& Hochberg, Y. (1995).
Controlling the false discovery rate: a practical and powerful approach
to multiple testing.
\emph{Journal of the Royal Statistical Society B}, 57(1), 289-300.

\bibitem[Christodoulou et al., 2019]{christodoulou2019systematic}
Christodoulou, E. et al. (2019).
A systematic review shows no performance benefit of machine learning
over logistic regression for clinical prediction models.
\emph{Journal of Clinical Epidemiology}, 110, 12-22.

\bibitem[Cohen, 1988]{cohen1988statistical}
Cohen, J. (1988).
\emph{Statistical Power Analysis for the Behavioral Sciences}, 2e ed.
Lawrence Erlbaum.

\bibitem[Collins et al., 2024]{collins2024tripodai}
Collins, G. S. et al. (2024).
TRIPOD+AI statement: updated guidance for reporting clinical prediction
models that use regression or machine learning methods.
\emph{BMJ}, 385, e078378.

\bibitem[Obermeyer et al., 2019]{obermeyer2019dissecting}
Obermeyer, Z., Powers, B., Vogeli, C., \& Mullainathan, S. (2019).
Dissecting racial bias in an algorithm used to manage the health of
populations.
\emph{Science}, 366(6464), 447-453.

\bibitem[Rubin, 1976]{rubin1976inference}
Rubin, D. B. (1976).
Inference and missing data.
\emph{Biometrika}, 63(3), 581-592.

\bibitem[Ruxton, 2006]{ruxton2006unequal}
Ruxton, G. D. (2006).
The unequal variance t-test is an underused alternative to Student's
t-test and the Mann-Whitney U test.
\emph{Behavioral Ecology}, 17(4), 688-690.

\bibitem[Tufte, 2001]{tufte2001}
Tufte, E. R. (2001).
\emph{The Visual Display of Quantitative Information}, 2e ed.
Graphics Press.

\bibitem[Wasserstein \& Lazar, 2016]{wasserstein2016asa}
Wasserstein, R. L. \& Lazar, N. A. (2016).
The ASA statement on p-values: context, process, and purpose.
\emph{The American Statistician}, 70(2), 129-133.

\bibitem[Weiskopf \& Weng, 2013]{weiskopf2013methods}
Weiskopf, N. G. \& Weng, C. (2013).
Methods and dimensions of electronic health record data quality
assessment.
\emph{Journal of the American Medical Informatics Association}, 20(1),
144-151.

\bibitem[Wong, 2011]{wong2011points}
Wong, B. (2011).
Color blindness.
\emph{Nature Methods}, 8(6), 441.

\end{thebibliography}

\end{document}
